\scanpage{33}\mspage{25}
% @ title, first part | text | Complexe cubique à balise marquée
\uline{Graphe cubique à balise marquée}
% @ title, first sign | diagram | triangle and stem, D₄
\diagram{\begin{tikzcd}[ampersand replacement=\&,execute at end picture={\draw(0,1)--(.66,1)--(.32,.54)--cycle (.32,.54)--(.32,-.07);\node[right] at(.9,1){$D_4$};}]\vphantom{A}\end{tikzcd}}
% @ title, second part | text | ou à paire de triangles
\uline{ou à paire} \deleted{\unsure{bitriangle épinglé}} \inserted{de triangles incidents}
% @ title, second sign | diagram | two triangles meeting at a vertex
\diagram{\begin{tikzcd}[ampersand replacement=\&,execute at end picture={\draw(0,1)--(.60,1)--(.30,.57)--cycle (.30,.57)--(.08,.17)--(.55,.17)--cycle;}]\vphantom{A}\end{tikzcd}}
marquée

% @ upper third, first case | text | Cas
Cas
% @ beside Cas | diagram | v,w ; u ; pendant segment
\diagram{\begin{tikzcd}[ampersand replacement=\&,execute at end picture={\draw[thick](0,1)node[above left]{$v$}--(.66,1)node[above right]{$w$}--(.36,.5)node[left]{$u$}--cycle; \draw(.36,.5)--(.36,-.10); \node[above]at(.33,1){$\unsure{\alpha|a}$};}]\vphantom{A}\end{tikzcd}}
, on a une arête marquée $\{v,w\}$, i.e. triangle $\{\deleted{\unsure{t}}u,v,w\}$ à sommet marqué, correspondant à un torseur sous \inserted{un} $\mathbb F_2^{P_u}$. Il faut en plus se donner un él. de $\widetilde P_u$, i.e. un élément de $P_u$ et une trivialisation du $\mathbb F_2$-torseur au dessus -- il reste donc un ens à 3 élmts $R_u=P_u-$

% @ mid-left, linked triangles | diagram | v,w,u,x,y ; dotted sides to y
\diagram{\begin{tikzcd}[ampersand replacement=\&,execute at end picture={
\draw(0,1.55)node[above left]{$v$}--(.79,1.52)node[above right]{$w$}--(.34,.81)node[left]{$u$}--cycle;
\draw(.34,.81)--(.35,-.13)node[below left]{$x$};\draw[dotted](.35,-.13)--(1.02,-.13)node[below right]{$y$}--(.34,.81);
}]\vphantom{A}\end{tikzcd}}
% @ middle to lower half | text | l'ens des 3 triangles
(l'ens des 3 triangles de sommet $u$ distincts de $\{u,v,w\}=t$ et $\{u,xy\}=s$) et un rev. d'ordre 2 de ce dernier, i.e. un torseur sous $\mathbb F_2^{R_u}$. La catégorie des triples $(C,\text{balise})$ est donc \unsure{décrite|discrète}, les autom d'un objet sont un groupe isom. à $\widehat W_{D_3}$ groupe de Weyl élargi (= tous les autom) d'un syst. de racines de type $D_3$,
\margin{balise épinglée donne un groupe d'autom de $(C,\text{bal. ép.})$ qui est $W_{D_3}$ (et non plus $\widehat W_{D_3}$)}% slanted
\inserted{il est de cardinal $3!\,2^3=\uline{48}$} (L'ens $\{v,w\}$ se récupère comme le produit \inserted{de Baer} des 3 \inserted{$\mathbb F_2$-}torseurs indexés par $R_u$ \ldots).

% @ lower third, roots | text | Les racines en question
Les racines en question sont les $\alpha-\beta$, où $\alpha,\beta\in E(\unsure{w|u})^*-\{x,y\}$ sont distincts et non liés.

% @ foot, count | text | Le nb des balises
Le nb des balises de $C$ est
% @ foot, product and labels | display | 45 . 3 . 8 = card W / 48
\[
\underbrace{45}_{\text{triangles }t}.\underbrace{3}_{\substack{\text{sommet }u\\\text{de }t}}.\underbrace{8}_{\in E(u)^*},\quad=\operatorname{card}W_{E_6}:48\qquad\text{OK.}/\ \deleted{\text{\unsure{Le nb des balises}}}
\]
% archive sheet number: 32
